\uuid{1mBx}
\titre{Coordonnées}
\auteur{bodin ; hanani ; mzari}
\organisation{exo7}
\licence{CC BY-NC-SA 4.0}
\datecreate{2020-01-20}
\niveau{L1}
\module{Géométrie}
\chapitre{Géométrie du plan}
\sousChapitre{Distance, norme, produit scalaire}
\difficulte{1}
\theme{coordonnées}
\exo7qcmid{lille-1.geometrie-plane.005}

\contenu{
\begin{question}
\qclassification{140.01,140.02}
\qtags{motcle=coordonnées}
Dans un repère orthonormé direct, on considère le point $A(1,1)$. Quelles sont les coordonnées polaires $(r,\theta)$ de $A$ ?
\begin{answers}
    \bad{$r=1$ et $\theta =1$}
    \bad{$r=2$ et $\theta =0$}
    \good{$r=\sqrt{2}$ et $\displaystyle \theta =\frac{\pi}{4}+2k\pi$, $k\in \Zz$}
    \bad{$r=\sqrt{2}$ et $\theta =0+2k\pi$, $k\in \Zz$}
\end{answers}
\begin{explanations}
D'abord, $r=\sqrt{1^2+1^2}=\sqrt{2}$ et $\theta $ est solution du système :
$$\left\{\begin{array}{l}\displaystyle \cos \theta =\frac{1}{\sqrt{2}}\\ \\ \displaystyle \sin \theta =\frac{1}{\sqrt{2}}.\end{array}\right.$$
\end{explanations}
\end{question}
}
